\documentclass[10pt,a4paper]{amsart}
\usepackage[margin=19mm]{geometry}
\usepackage{amsmath,amssymb,mathtools}
\usepackage[scr=rsfs,cal=euler]{mathalfa}
\usepackage{xfrac}
\usepackage[unicode,hidelinks,bookmarks=false]{hyperref}
\newcommand{\ie}{i.e.\ }
\newcommand{\cf}{cf.\ }
\newcommand{\resp}{respectively\ }
\newcommand{\Subsection}{\subsection}
\providecommand{\nobreakdash}{\nobreak\hbox{-}}
\setlength{\emergencystretch}{2em}
\multlinegap0pt
\usepackage{amssymb}
\newtheorem{theorem}{Theorem}
\title{Hindman and Owings-like theorems \\ without the Axiom of Choice}
\author{José Alberto Guzmán-Vega, David José Fernández-Bretón, Eliseo Sarmiento-Rosales}
\date{Source: arXiv:2603.27163v4 (CC BY 4.0)}
\begin{document}
\maketitle

\noindent\emph{Source and licence.} Hindman and Owings-like theorems \\ without the Axiom of Choice. Version arXiv:2603.27163v4 (CC BY 4.0), distributed by arXiv under the Creative Commons Attribution 4.0 International licence, http://creativecommons.org/licenses/by/4.0/, as recorded in the arXiv licence metadata for this exact version and shown on the abstract page https://arxiv.org/abs/2603.27163v4. Full licence notice: \texttt{licenses/arxiv-2603.27163-CC-BY-4.0.txt}.\par\smallskip

\noindent\textit{Editorial scope.} Counted unit: the article's theorem fin-fin-hindman (source0.tex lines 184--186). The article's complete original proof of it is delivered with it (source0.tex lines 187--211, one \texttt{proof} environment of 25 lines). No mathematics is written, repaired or re-derived: the statement and the proof are the article's own lines, in the article's own notation, and no retained line is substituted or re-flowed.  No further source-local context is needed: every cross-reference of every retained line resolves inside the two delivered spans.  Nothing was omitted from any retained span, so the package carries no omission record.  No retained line contains a citation, so the wrapper carries no bibliography and no bibliography entry.  No retained line contains a figure environment, an image-inclusion command or drawing code, so no image is delivered and the package declares no asset dependency.  Editorial formatting only, in addition: an AMS \texttt{amsart} preamble that restates the article's own declarations and macros used by the retained lines, namely the environments \texttt{theorem} on separate counters, together with the standard packages those lines need.  The retained environments print in this document as Theorem 1 in order of appearance, whereas the article prints the same items with the numbering of its own full document.  The complete native source of the article as distributed in the arXiv e-print for this exact version is bundled unchanged as \texttt{sources/197326/source0.tex}.
\begin{theorem}\label{fin-fin-hindman}
    For all finite $\kappa,\theta$ and $L\geq 1$, there exists $S\in\mathbb{N}$ such that, for every $L$-left cancellative semigroup $(G,+)$ with $|G|\geq S$, $G\to (\kappa)_{\theta}^{\operatorname{FS}}$.
\end{theorem}

\begin{proof}
    First, use a compactness argument with Hindman's original theorem on $\mathbb N$ to obtain a finitary version of the finite-unions version of Hindman's theorem: given finite $\theta$ and $\kappa$, there exists some natural number $F$ such that, for every colouring of $\wp(F)\setminus\{\varnothing\}$ in $\theta$ colours, there is a block sequence $B\subseteq\wp(F)\setminus\{\varnothing\}$ of size $\kappa$ such that $\operatorname{FU}(B)$ is monochromatic. Similarly (or by mapping each finite $a\subseteq\mathbb N$ to the number $\sum_{i\in a}2^i$), we obtain a finitary version of the finite-sums version of Hindman's theorem: for each finite $\theta$ and $\kappa$, there is a number $R$ (it suffices to take $R=2^{F}-1$ where $F$ is as in the previous statement) such that, for every colouring of $\{1,\cdots,R\}$ in $\theta$ colours, there is $M\subseteq R$ of size $\kappa$ such that $\operatorname{FS}(M)$ is monochromatic.
    
    Now, given finite $\kappa, \theta$ and $L\geq1$, consider the numbers $F, R$ as described in the previous paragraph, and fix $S:=\max\{R+1, 4^{F}(L+1)\}$.
    Consider any $L$-left cancellative semigroup $G$ with $|G|\geq S$ and $c:G\longrightarrow\theta$. Letting $\langle g\rangle$ denote the monogenic subsemigroup generated by $g\in G$, we observe that the following two cases are exhaustive:
    \begin{description}
        \item[$\exists g\in G:|\langle g\rangle|>R$]
            In this case, the set $\{1, 2, \cdots, R\}$ is in bijection with $\{g, 2g, \cdots, Rg\}$; furthermore, addition is preserved whenever the sum stays within the set. Thus, it suffices to consider the induced colouring on $\{1,\cdots,R\}$ given by the restriction $c\upharpoonright\{g,2g,\cdots,Rg\}$, and apply the finite version of Hindman's theorem.
        \item[$\forall g\in G:|\langle g\rangle|\leq R$] 
            For all $n<F$, recursively define:
            \begin{align*}
                E_{n}:=&\operatorname{FS}(h_{j}\mid j<n),\\
                D_{n}:=&\{g\in G\mid\exists e\in E_{n}\cup\{0\}\;\exists f\in E_{n}\;(e+g=f)\}.
            \end{align*}
            The $L$-left cancellativity of $G$ implies that:
            \begin{align*}\;
                |D_{n}\cup E_{n}|\leq
                |E_{n}\cup\{0\}||E_n|L+|E_n|<
                4^{n}L+2^{n}<
                4^{n}(L+1),
            \end{align*}
            therefore there is at least one $h_{n}\in G\setminus(D_{n}\cup E_{n})$ whenever $n<F$.
            Now consider the map $d:\wp(F)\setminus\{\varnothing\}\to G$ given by $d(f)=\sum_{i\in f}h_{i}$ (ordered by index) and apply the finite finite-unions theorem to $c\circ d$, so that there exists a block sequence $B=(b_{i}\mid i<\kappa)$ in $\wp(F)\setminus\{\varnothing\}$ such that $\operatorname{FU}(B)$ is $c\circ d$-monochromatic. This readily implies that, if we let $M:=(d(b_{i})\mid i<\kappa)$, $\operatorname{FS}(M)$ is $c$-monochromatic. \qedhere
    \end{description}
\end{proof}

\end{document}
